%------------------------------------------------------------------------------
% CHAPTER 1: INTRODUCTION
%------------------------------------------------------------------------------
\chapter{Introduction}
\label{ch:introduction}

\section{Purpose}
\label{sec:purpose}

The Best Bike Paths (\texttt{BBP}) system addresses the recurring need for Bike Riders to discover new routes and track their cycling activities. Registered users can record their rides either manually, by entering street names, or automatically, leveraging their device sensors.

All users, regardless of their registration status, can search for paths between any given origin and destination point, visualizing them on the map. Search results are ranked based on a score that integrates the path's status and its efficiency, measured as trip time. Additionally, the system offers statistics for each path to any user, including trip time and, when available, it also provides the current weather conditions along the route.

A registered user can access his previously recorded and validated trips and look at its performance, speed, and weather conditions during the recording.

The system aggregates identical paths submitted by multiple users, using metrics like freshness and status to combine the information. As part of this merging process, the final trip time is calculated by averaging the individual trip times of the paths being aggregated.

\subsection{Goals}
\label{subsec:goals}

\begin{itemize}
    \item \textbf{[G0]:} A user wants to signup first and then login to the application.
    \item \textbf{[G1]:} A registered user aims to automatically record their cycling trips using theirs mobile device's sensors. 
    \item \textbf{[G2]:} A registered user aims to manually record a bike trip by inserting one street at a time without physical cycling.
    \item \textbf{[G3]:} A registered user wants to access their recorded bike trips history.
    \item \textbf{[G4]:} A registered user aims to view trip on map and consult its statistics, including average speed, trip time and meteorological conditions (at recording time).
    \item \textbf{[G5]:} A registered user wants to share information about a bike trip they have recorded, including their status and the presence of obstacles, making this information available to the community (all users).
    \item \textbf{[G6]:} A User wants to search for bike paths given an origin and a destination point and visualize results on a map.
    \item \textbf{[G7]:} A User wants to retrieve up-to-date information about bike paths, their status, and current weather conditions.
    \item \textbf{[G8]:} A registered user wants to delete one of their previously recorded trips.
\end{itemize}

\section{Scope}
\label{sec:scope}

\subsubsection{Main Features}

\paragraph{Trip Recording}
\begin{itemize}
    \item Manual entry of bike paths with street-level details and path status
    \item Automatic tracking of rides using mobile device sensors, including distance, average speed and duration
    \item Ensuring continuous data collection and storage of the current trip even in areas with no connectivity, to be synchronized once the connection is restored
\end{itemize}

\paragraph{Path Discovery and Information}
\begin{itemize}
    \item Search for bike paths between specified locations
    \item Visualization of paths on map
    \item Display of path characteristics (total distance, trip time)
    \item Provision of current weather information for a selected route
\end{itemize}

\paragraph{Up to date Information}
\begin{itemize}
    \item Aggregation of path data from multiple users
\end{itemize}

\paragraph{Personal Analytics}
\begin{itemize}
    \item Storage and retrieval of a registered user cycling history
\end{itemize}



\subsection{World Phenomena}

\begin{enumerate}
    \item A Registered User wants to go for a bike ride and record the trip.
    \item A User wants to get all paths available given origin and destination point.
    \item A registered user wants to get all his private trips.
    \item A User wants to know information about a specific bike path, such as total track distance, average speed and status.
    \item A registered user wants to record a bike trip in manual mode.
    \item A registered user wants to record a bike trip in automatic mode.
    \item A registered user wants to publish one of his private trips.
    \item A user wants to validate obstacle on the trip (within a certain time interval from its recording).
    \item A mobile device with sensors (GPS, accelerometer, gyroscope) collects GPS and motion data while the user is biking.
    \item External weather data services provide information about environmental conditions.
\end{enumerate}

\subsection{Shared Phenomena}

\subsubsection{World Controlled}

\begin{enumerate}
    \item A user signs up to the system.
    \item A registered user starts recording their trip (automated mode).
    \item A registered user stops recording their trip (automated mode).
    \item A registered user inserts their bike trip data, specifying name and status of each street (manual mode).
    \item A registered user accesses their private bike trips.
    \item A registered user publishes their trip.
    \item A user searches for all paths available from an origin and a destination point.
    \item A user consults the list of paths on the map.
    \item A user selects one of the proposed paths.
    \item A registered user confirms or denies the presence of a pothole (automatically detected by BBP).
\end{enumerate}

\subsubsection{Machine Controlled}

\begin{enumerate}
    \item The system shows bike trip data acquired from mobile devices of the registered user (automated mode).
    \item The system asks the registered user to confirm or correct the presence of a pothole.
    \item The system shows on the map the bike path(s) between given origin and destination points.
    \item The system rank the paths based on their score (involving path status and effectiveness)
    \item The system displays consolidated path information, merging data provided by different users
\end{enumerate}

\section{Definitions, Acronyms, Abbreviations}
\label{sec:definitions}

\subsection{Definitions}
\begin{itemize}
    \item \textbf{Registered user}: User that is registered in the platform.
    \item \textbf{Guest user}: User that is not registered in the platform.
    \item \textbf{User}: Term used to indicate both registered and guest users.
    \item \textbf{Path}: Public route connecting two geographical points, obtained from the merging process of multiple Trips sharing identical GPS coordinates and cumulative distance metrics. This designation indicates either the presence of dedicated cycling infrastructure or low-traffic roadways with speed limits compatible with typical cycling velocities.
    \item \textbf{Trip}: A user-specific ordered sequence of GPS coordinates or street identifiers, accompanied by associated meteorological data, recorded by a registered user.
    \item \textbf{Merging process}: A computational process through which the \texttt{BBP} system aggregates trip data, potentially coming from multiple users, to generate public paths. It also represents the process of updating an already-existing public path with recently published trips.
    \item \textbf{Anomaly}: Pothole or obstacle found while recording the trip and identified by the sensors data analysis.
    \item \textbf{Trip status}: The maintenance condition detected or validated for a user trip, based on sensor data analysis (potholes) or manual insertion.
    \item \textbf{Path status}: The aggregated maintenance condition of a bike path, computed by the \texttt{BBP} system whenever a new trip is published and merged. The calculation considers the status of all merged trips, weighing factors such as data freshness and consensus among multiple users' reports (e.g., if two trips report "optimal" and one reports "requires maintenance" within the same timeframe, the path is considered "optimal").
    \item \textbf{Path score}: A quality metric used to rank and visualize bike paths, calculated based on the path's status and its effectiveness in connecting origin to destination (measured through average trip duration).
    \item \textbf{Corresponding Path (to a Trip)}: The path that shares the same trajectory (ordered set of GPS coordinates) as a given trip.
    \item \textbf{Trip Merging Timeframe}: The time window within which trip reports are considered for consensus-based status aggregation during the path merging process.
\end{itemize}

\subsection{Acronyms}
\begin{itemize}
    \item \textbf{API}: Application Programming Interface
    \item \textbf{BBP}: Best Bike Paths
    \item \textbf{GPS}: Global Positioning System
    \item \textbf{GDPR}: General Data Protection Regulation
\end{itemize}

\section{Reference Documents}
\label{sec:references}
\begin{itemize}
    \item Project assignment document
    \item Software Engineering 2 Course Material - A.Y. 2025/2026
\end{itemize}

\section{Document Structure}
\label{sec:structure}
This document is organized as follows:
\begin{itemize}
    \item \hyperref[ch:introduction]{\color{bluepoli}C1} \color{black} \textbf{Introduction}: We introduce the goals of the project, purposes, and an analysis on world and shared phenomena; we also list abbreviations and definitions, needed to fully comprehend the document
    \item \hyperref[ch:overall]{\color{bluepoli}C2} \color{black} \textbf{Overall Description}: We provide an overview of the system through possible scenarios and explain the functions of the various Users. To complete the description, we add the assumptions and dependencies at the base of the system.
    \item \hyperref[ch:requirements]{\color{bluepoli}C3} \color{black} \textbf{Specific Requirements}: Provides a detailed analysis of both functional and non-functional system requirements. In this chapter, Use Cases Diagrams are also provided, to show the possible system interactions. Lastly, we add traceability matrixes to help understand relations between goals, requirements, domain assumptions and use cases.
    \item \hyperref[ch:alloy]{\color{bluepoli}C4} \color{black} \textbf{Alloy Formal Analysis}: Provides a formal analysis of the system to be using Alloy 6
    \item \hyperref[ch:effort]{\color{bluepoli}C5} \color{black} \textbf{Effort Spent}: Describes and shows how the work load has been distributed between each team members
    \item \color{bluepoli} C6 \color{black}\textbf{References}: Provides a list of the references used to draw up this document.
\end{itemize}


